% Drawing 152/01-8 -- folder 152, batch 01, archive page 1 (PDF page 2).
% Transcription: au milieu de la moitié gauche, au crayon : un arbre de flèches. Une tige fléchée descend jusqu'à un sommet, d'où partent une flèche vers le haut à droite, une branche fléchée vers le bas à gauche et une branche horizontale repassée d'un trait épais, portant trois pointes de flèche ; du bout de celle-ci, une branche fléchée monte à droite jusqu'à un second sommet et se prolonge vers $s_3$, une autre descend à droite ; du second sommet part vers la droite un trait qui se divise en deux.
% Kind: pencil tree of arrows. Method 2 (schematic redraw from the facsimile). Status: ok.
% Notes: The stem coming down to the first vertex with three arrowheads on it, the short arrow to the upper right, the branch to the lower left with two arrowheads, the heavy horizontal branch with three heavy arrowheads (drawn with a double stroke on the page, here one thick line), the branch rising to the second vertex with two arrowheads and its continuation up to the chain of 01-6 (it ends under s_3), the branch going down to the right with a small tick across it, and the faint stroke leaving the second vertex to the right and forking in two. Every branch is drawn from the computed vertex. Overlaid on the 400 dpi crop: the strokes fall on the page's. Not drawn: the ink frame of the list at the right, which the fork runs over. Redrawn 2026-11-12 (issue 128): to check again.
% Coordinates: 1 unit = 1 mm = 1/100 of the crop width (bbox in page fractions: [0.25, 0.12, 0.56, 0.27]).
% Redrawn by Claude (claude-fable-5-2), issue #128, 2026-11-12; not checked by a human; not compiled (no TeX in the session).
% Rights: reproduces part of the work -- behind RENDER_TRANSCRIPTS, never published by itself.
\documentclass[tikz,border=3pt]{standalone}
\usepackage{amsmath,amssymb}
\usetikzlibrary{patterns,arrows.meta}
\tikzset{pencil/.style={black!45}, struck/.style={black!35}, mpt/.style={circle,draw,inner sep=1.3pt}, pt/.style={circle,fill,inner sep=0.9pt}, opt/.style={circle,draw,fill=white,inner sep=1.1pt}}
\begin{document}
\begin{tikzpicture}[x=1mm,y=1mm,line cap=round,line join=round,font=\small]
\draw[-{Stealth[length=1.8mm]}, pencil] (20.6,-3) -- (20.5,-10.5);
\draw[-{Stealth[length=1.8mm]}, pencil] (20.5,-10.5) -- (20.1,-22.5);
\draw[-{Stealth[length=1.8mm]}, pencil] (20.1,-22.5) -- (19.5,-39.2);
\draw[-{Stealth[length=1.8mm]}, pencil] (19.5,-39.2) -- (30.6,-29.8);
\draw[-{Stealth[length=1.8mm]}, very thick] (19.5,-39.2) -- (30,-38.2);
\draw[-{Stealth[length=1.8mm]}, very thick] (30,-38.2) -- (37,-37.5);
\draw[-{Stealth[length=1.8mm]}, very thick] (37,-37.5) -- (44.6,-36.6);
\draw[-{Stealth[length=1.8mm]}, pencil] (44.6,-36.6) -- (54,-28);
\draw[-{Stealth[length=1.8mm]}, pencil] (54,-28) -- (59.3,-24.2);
\draw[pencil] (59.3,-24.2) -- (64.5,-20.5);
\draw[pencil] (64.5,-20.5) -- (74.2,-2.5);
\draw[-{Stealth[length=1.8mm]}, pencil] (44.6,-36.6) -- (59.6,-48);
\draw[pencil] (51.8,-41.6) -- (52.6,-43.2);
\draw[-{Stealth[length=1.8mm]}, pencil] (19.5,-39.2) -- (11.5,-54);
\draw[-{Stealth[length=1.8mm]}, pencil] (11.5,-54) -- (5,-62.3);
\draw[pencil] (64.5,-20.5) -- (85,-23);
\draw[pencil] (85,-23) -- (91,-11.5);
\draw[pencil] (85,-23) -- (94,-34.5);
\end{tikzpicture}
\end{document}
